\section{The Codex as Flesh, Stone, and Steel}

A traditional Runekeeper's Codex is a book. Vellum pages, leather binding, ink that smells of lampblack and blood. It is portable, concealable, and replaceable. It is also fragile. Fire destroys it. Water ruins it. A thief can steal it.

A Runic Warrior's Codex is not a book.

It is a gorget engraved with vows. A shield marked with the names of thresholds guarded. A staff carved with seasonal runes. A sword with Rites etched along the fuller. A body marked by tattoo, brand, or scar.

The Codex is not carried. It is worn. It is held. It is the body itself.

\subsection{The Three Forms}

A Runic Warrior's Codex takes one of three forms, each with advantages and risks.

\subsubsection{Flesh (Tattoos, Brands, Scarification)}

The Rites are written on the Runekeeper's skin. Needle and ash, blade and ink, fire and clay. The Codex cannot be stolen. It cannot be confiscated at a city gate. It is always present.

But it is also always visible. Tattoos curl up the neck. Brands mark the back of the hand. Scars tell stories that cannot be hidden.

\textbf{Advantages:}
\begin{itemize}
\item Cannot be stolen or confiscated.
\item Always accessible (no action to draw or retrieve).
\item Counts as a visible mark of covenant (may intimidate or impress).
\end{itemize}

\textbf{Risks:}
\item Cannot be hidden without clothing that covers every inch.
\item Damage to the skin (burns, deep wounds) may damage the Codex.
\item Some cultures view flesh-writing as heretical or criminal.

\subsubsection{Stone (Carved Objects, Shields, Weapons, Armor)}

The Rites are carved, etched, or engraved onto a physical object. A gorget, a shield, a staff, a sword. The Codex is durable. It can be used as a tool or weapon. It can be displayed as a symbol of authority.

But it can be lost. It can be stolen. It can be broken.

\textbf{Advantages:}
\begin{itemize}
\item Can serve as armor (gorget, shield) or weapon (staff, sword).
\item Durable—resistant to fire, water, and mundane damage.
\item Visible symbol of covenant (may grant social advantage in certain contexts).
\end{itemize}

\textbf{Risks:}
\item Can be stolen, confiscated, or lost.
\item Can be damaged or destroyed (requires repair or replacement).
\item Heavy or conspicuous (may be difficult to conceal).

\subsubsection{Steel (Sentient Weapons)}

The Rites are engraved on a sentient blade or object. The weapon speaks. It remembers. It serves as both Codex and Thiasos, though it cannot act independently.

This is the rarest form. Sentient weapons are not made; they are found, inherited, or granted by a Patron. They have personalities, opinions, and sometimes demands.

\textbf{Advantages:}
\begin{itemize}
\item Functions as both Codex and Thiasos (counts as a witness for Rites).
\item Cannot be stolen without the Runekeeper's awareness (the weapon may resist or cry out).
\item Provides counsel, warnings, and memory (GM may use the weapon to deliver hints).
\end{itemize}

\textbf{Risks:}
\item The weapon has its own personality and may refuse to cooperate.
\item Cannot be easily replaced or repaired.
\item Attracts attention—sentient weapons are rare and coveted.

\subsection{Engraving Rites}

A Runic Warrior adds new Rites to their Codex through physical engraving. This is not a scholarly exercise. It requires a forge, a needle, a chisel, or a branding iron. It requires pain.

\textbf{Mechanics:}

Adding a new Rite to a Codex requires:
\begin{itemize}
\item Downtime (one day per Tier of the Rite).
\item A suitable workspace (forge, studio, quiet room).
\item Materials appropriate to the Codex form (ink, ash, metal, leather, wood).
\item A Body + Craft test (DV = 3 + Rite Tier).
\end{itemize}

On a success, the Rite is engraved. On a failure, the materials are ruined, and the Runekeeper marks 1 Fatigue (the pain of the failed attempt). On a critical failure, the Codex is damaged (becomes Compromised) and must be repaired before new Rites can be added.

\subsection{Repairing a Damaged Codex}

A Codex can be damaged. Fire, water, blade, or simple wear.

\begin{itemize}
\item A \textbf{Flesh} Codex can be repaired through healing (rest, Medicine, or magical healing). Damaged tattoos fade; scars may be re-cut.
\item A \textbf{Stone} Codex requires physical repair (Craft test, DV 4, downtime). A broken shield can be reforged. A cracked staff can be bound with wire.
\item A \textbf{Steel} Codex (sentient weapon) requires the Runekeeper to negotiate with the weapon itself. The weapon may demand a service, a memory, or a promise before it allows repairs.
\end{itemize}

If a Codex is destroyed beyond repair, the Runekeeper loses all Rites engraved on it. They must acquire or craft a new Codex and re-engrave their Rites—a costly and time-consuming process.

\subsection{The Codex as Witness}

A Codex is not just a record. It is a witness. When a Runic Warrior invokes a Rite, their Codex sees. It remembers. This counts as a witness for Rites that require one (improving Position or reducing DV).

A sentient Codex (Steel form) can also speak. It may offer testimony, confirm an oath, or testify to a covenant. In a hedge-court, a sentient blade's word may carry as much weight as a living witness.

\subsection{The Codex and Obligation}

When a Runic Warrior marks Obligation, the cost is visible. A scar appears. A tattoo darkens. A crack spreads across the shield. The Codex records the debt in the same way the body records a wound.

This visibility is a risk. An Inquisitor can see the Obligation. A Chain-Lantern can count the scars. But it is also a warning: the Runekeeper cannot pretend to be neutral. Their covenant is written on them.

The thread held. The water carried. That is not a Rite. That is grace.
